\chapter{研究点一 多品种证券担保品估值及价差风险计量技术}

\section{引言}

【说明本研究点解决的问题及其与总体目标的关系。】

\section{多品种担保品估值与风险计量相关研究}

【评述与本研究点方法设计直接相关的工作。】

\section{M1：全部 A 股估值及价差风险模型}
\label{sec:m1-equity}

\subsection{模型定位与计量口径}

M1 面向全部 A 股建立股票价格、波动率与持有期下行风险的联合计量模型。模型以估值时点可获得的复权价格、成交量、换手率、停牌和涨跌停状态等信息为输入，在统一复权口径下生成未来价值的条件分布。分红、送转等公司行为必须按照事前确定的规则进入价格或现金流口径，不能同时计入复权价格和期内现金流。对资产 \(i\)、估值日 \(t\) 和持有期 \(h\)，M1 同时输出 \(V_{i,t+h}\) 的点预测、区间和分位数，并由同一价值路径得到
\begin{equation}
  \Delta V_{i,t,h}=V_{i,t+h}-V_{i,t},\qquad
  L_{i,t,h}=V_{i,t}-V_{i,t+h},
\end{equation}
以及波动率、VaR、ES、Amihud 非流动性指标、换手率、压力损失和预测区间宽度等结果\cite{amihud_2002,artzner_1999}。

“全部 A 股”是覆盖要求，而不是抽取若干代表性股票后的外推结论。为此，模型运行需逐日生成全量证券目录、合格样本目录和排除台账，分别核算数据资格率、模型运行覆盖率和可回测率；新股、长期停牌、退市整理期和公司行为异常样本应按事前规则处理并保留原因码。

\subsection{GBSV-v4.1 条件价值路径}

M1 的研究候选保持可识别的 GBSV 模型家族。其结构由四个相互衔接的部分组成：首先使用仅依赖时点 \(t\) 及以前信息的因果滤波形成局部波动状态；其次以 GIGa/广义 Boltzmann 边际刻画正值波动率的厚尾横截面；再次以 Gaussian copula 表示相邻波动状态的动态持续性；最后通过带 Itô 修正的价格方程把波动率路径转化为价值路径。第 \(m\) 条模拟路径可写为
\begin{equation}
  V_{t+h}^{(m)}=V_t\exp\!\left\{
  \sum_{j=1}^{h}\left[
  \mu_{t+j}^{(m)}-\frac{1}{2}\bigl(\sigma_{t+j}^{(m)}\bigr)^2
  +\sigma_{t+j}^{(m)}z_{t+j}^{(m)}
  \right]\right\}.
\end{equation}
当前 v4.1 主候选冻结为零漂移、无杠杆版本，目的是避免在有限样本中用不稳定的漂移或杠杆参数放大路径均值。分层状态、市场因子和神经参数头可作为后续扩展，但不得移除 GBSV 的波动边际、动态持续性与条件路径核心，也不得把非 GBSV 学习器静默替换为 M1。

每条价值路径均转换为 T+1、T+5 和 T+10 的价值差、正损失及风险分位数。对置信水平 \(q\)，VaR 与 ES 由同一路径损失样本计算：
\begin{equation}
  \operatorname{VaR}_{q}=Q_q(L),\qquad
  \operatorname{ES}_{q}=\operatorname{E}\!\left[L\mid L\geq Q_q(L)\right].
\end{equation}
这一设计使价格预测、波动率预测和尾部风险评价共享相同的信息集和情景路径，避免分别训练的输出在方向、期限或损失口径上相互矛盾。

\subsection{基准、验证与证据边界}

M1 的朴素基准包括上一期价值、历史均值和历史波动率；经典基准包括 GBM、GARCH 和历史模拟\cite{engle_1982}；TCN、Transformer 或 boosted panel 等方法仅作为强挑战基准。所有方法必须使用相同的合格样本、预测起点和信息集，分别比较价值 MAE/MASE、波动率 QLIKE、分位数 pinball loss、区间覆盖及 VaR 返回检验。

现有证据必须分层报告。旧版 GBSV-v2 未获得稳定科学增益；v3 虽然工程运行成功，但由于未加入 \(-\frac{1}{2}\sigma^2\) 修正，厚尾波动混合在价格指数化时产生极端伪均值，因此判定为科学失败并永久保留为负向证据。修正后的 v4.1 在排除开发股票的 96 只股票有界横截面 holdout 上通过预注册数值稳定性、价值 MASE 非劣和聚合 QLIKE 门禁，但 GBM 在 T+5、T+10 点预测上仍更强。故当前状态仅为研究候选，不能据此声称已经实现全 A 股覆盖、时间外稳定性或合同准确率通过。

\section{M2：公募债券估值及价差风险模型}
\label{sec:m2-bond}

\subsection{模型定位与计量口径}

M2 面向公募债券建立现金流一致的全价估值和持有期总回报损失模型。估值主口径采用 dirty value；若数据源只提供净价，则必须显式加回应计利息并与票息、本金和提前偿付记录对账。债券在持有期内通常存在现金流，仅比较期初和期末净价会把票息误记为损失，因此正式风险口径为
\begin{equation}
  \operatorname{PnL}_{t,h}=V_{t+h}^{\mathrm{dirty}}-V_t^{\mathrm{dirty}}+CF_{(t,t+h]},
  \qquad L_{t,h}^{\mathrm{TR}}=-\operatorname{PnL}_{t,h},
\end{equation}
其中 \(CF_{(t,t+h]}\) 包括持有期内实际发生的票息、本金、提前偿付和违约回收。模型同时保留不含现金流的价值变化作为诊断字段，但不得与总回报损失混用。

\subsection{现金流、曲线与条件价值路径}

债券价值主干由合约现金流、无风险收益率曲线、信用利差和回收假设共同构成。对尚未支付的第 \(k\) 笔现金流，时点 \(t\) 的理论全价可概括为
\begin{equation}
  V_t^{\mathrm{dirty}}
  =\sum_{k:T_k>t}\operatorname{E}_t(CF_k)D_t(T_k)
  +R_t^{\mathrm{model}},
\end{equation}
其中 \(D_t(T_k)\) 为与期限和信用状态一致的折现因子，\(R_t^{\mathrm{model}}\) 为在严格时间切分下学习的定价残差。收益率曲线可采用 Nelson--Siegel 等参数化结构作为经典锚\cite{nelson_siegel_1987}；信用债则以强度、迁移或利差过程反映违约概率和风险溢价\cite{duffie_singleton_1999}。残差学习只能校正可解释主干，不能替代现金流对账或使用估值日之后的评级、成交和曲线信息。

未来情景同时推进基准曲线、信用利差和流动性状态，并在每条路径上重估剩余现金流。模型据此输出未来全价、总回报 PnL、VaR、ES、收益率变化、曲线利差、信用利差、久期、DV01、成交活跃度和压力损失。对浮息、含权、可转债及资产支持证券，应先按产品条款分层；不能用普通固定利率债券的单一现金流规则覆盖全部子类。

\subsection{基准、验证与证据边界}

M2 的朴素基准为上一期收益率、利差或全价；经典基准为现金流贴现、收益率曲线和 reduced-form 信用定价；强现代基准可采用 boosted residual 或时序/面板学习器。比较必须使用相同债券目录、现金流版本、估值日期和持有期，并分别核验价格误差、收益率误差、现金流对账差异、分布得分和 VaR/ES 返回检验。

当前债券数据已具备全价、应计利息、收益率、久期和 DV01 等字段，原则上可形成 T+1、T+5 和 T+10 标签；但票息、本金、提前偿付、违约事件与无风险曲线仍需逐项审计。因而 M2 现阶段可以开展工程 smoke 和小型时间外验证，尚不能把文件数量或局部定价成功率表述为全部公募债券覆盖，更不能在现金流未闭合时给出正式验收结论。

\section{M3：公募基金估值及价差风险模型}
\label{sec:m3-fund}

\subsection{模型定位与计量口径}

M3 面向公募基金建立净值、可交易价格和持有期损失的分层计量模型。不同基金的可观察价值并不相同：开放式基金通常以单位净值为主，上市基金还存在交易价格及折溢价，货币基金、债券基金、股票基金、指数基金和跨境基金的估值频率及风险来源亦有明显差异。因此，基金类型、主价值口径和可交易性必须在建模前冻结，不能把 NAV 与市场价格混成同一目标。

对采用总回报口径的基金，持有期 PnL 定义为
\begin{equation}
  \operatorname{PnL}_{t,h}
  =V_{t+h}-V_t+D_{(t,t+h]}+A_{(t,t+h]},
\end{equation}
其中 \(D_{(t,t+h]}\) 表示分红，\(A_{(t,t+h]}\) 表示按事前规则计量的申赎、折算或清盘调整。上市基金还需同时报告折溢价及其变化，用于解释交易价格与净值路径的偏离，但折溢价本身不替代跨产品统一的持有期价值损失。

\subsection{持仓穿透、因子复制与残差学习}

M3 以“可观察净值/价格路径+持仓穿透或市场因子复制+受限残差学习”为主干。对于持仓披露充分的基金，使用股票、债券、商品和现金仓位构造复制价值；对于持仓低频披露或存在滞后的基金，以基金类型、基准指数和收益因子建立状态空间或动态因子路径，并显式标记持仓信息的发布日期。条件价值可写为
\begin{equation}
  \widehat V_{t+h}
  =V_t\left(1+\boldsymbol{\beta}_{t}^{\mathsf T}
  \widehat{\boldsymbol r}_{t,t+h}\right)
  +\widehat\varepsilon_{t,h},
\end{equation}
其中 \(\boldsymbol{\beta}_{t}\) 为时点可得的持仓或因子暴露，
\(\widehat{\boldsymbol r}_{t,t+h}\) 为市场因子情景，
\(\widehat\varepsilon_{t,h}\) 为按严格时间切分估计的跟踪残差。该结构既保留资产来源的可解释性，也允许在不泄漏未来持仓的前提下刻画非线性跟踪误差。

每条情景输出未来 NAV 或交易价格、持有期总回报损失、VaR、ES、折溢价、跟踪误差、申赎压力、持仓流动性、集中度和区间宽度。若某项输入不可观察，应返回明确的不可观察状态，而不是以零值填充。

\subsection{基准、验证与证据边界}

M3 的朴素基准为上一期 NAV、交易价格或收益率；经典基准为基准指数、因子模型或持仓复制；强现代基准为 boosted 或 sequence residual model。估值与风险比较应在基金子类内进行，并统一估值频率、信息发布日期和分红调整。主要评价包括 NAV/价格 MAE 与 MASE、跟踪误差、折溢价误差、CRPS、pinball loss 及 VaR/ES 返回检验。

现有基金数据可以支持上市产品的价格、折溢价和成交特征分析，但上市与非上市基金的价值口径、日频 NAV 可得性以及分红、折算和清盘规则尚未统一。故当前工作重点是完成基金类型分层和价值口径审计；在此之前，局部可运行模型只能作为候选组件，不能代表全部公募基金的统一结论。

\section{M4：场内期货估值及价差风险模型}
\label{sec:m4-futures}

\subsection{模型定位与计量口径}

M4 面向场内期货建立结算价、期限结构和逐日盯市损益的联合模型。期货与现货资产不同，其风险不仅来自价格变化，还受合约乘数、头寸方向、保证金、到期和展期规则影响。对头寸数量 \(q\)、方向 \(s\in\{-1,1\}\) 和合约乘数 \(M\)，不跨展期的持有期累计盯市损益可写为
\begin{equation}
  \operatorname{PnL}_{t,h}=s\,qM\bigl(F_{t+h}-F_t\bigr),
  \qquad L_{t,h}=-\operatorname{PnL}_{t,h}.
\end{equation}
若持有期跨越到期或预设换月日，则必须按冻结的可交易规则加入平仓、移仓和费用现金流，不能把事后拼接的连续合约价格直接当作可实现损益。

\subsection{持有成本、期限结构与结算价路径}

M4 以持有成本关系和期限结构为经典锚。对适用的合约，理论远期价格可表示为
\begin{equation}
  F_{t,T}^{\mathrm{theory}}
  =S_t\exp\{(r_t+u_t-y_t)(T-t)\},
\end{equation}
其中 \(r_t\) 为资金成本，\(u_t\) 为仓储、保险或其他持有成本，\(y_t\) 为便利收益或等价持有收益。金融期货和商品期货的变量含义及约束不同，模型需按品种组估计基差、月间价差和期限结构状态，再对结算价或基差残差进行条件学习。

未来情景联合推进现货或指数因子、近远月价差、波动率、成交和持仓状态，并按合约日历执行逐日盯市、到期和展期。主要输出包括未来结算价、累计 MTM 损益、VaR、ES、基差、月间价差、展期收益、保证金压力、成交/持仓流动性和压力损失。风险输出应同时保留实际合约与研究用连续序列的标识，防止两类价格被混用。

\subsection{基准、验证与证据边界}

M4 的朴素基准为上一结算价或上一基差；经典基准为持有成本模型和历史模拟；强现代基准可采用期限结构 sequence model 或 boosted model。验证应按合约、品种和期限分层，在相同换月规则下比较结算价/基差误差、MTM 分布得分、压力期表现和 VaR/ES 返回检验。

现有期货数据覆盖结算价、收盘价、基差、持仓和成交等关键字段，具备建立真实标签的基础；但合约乘数、头寸方向、保证金、到期与连续合约展期规则仍需统一冻结。因而 M4 当前属于可推进但合同规则尚未闭合的模型，研究结果必须明确实际合约样本与连续合约样本的边界。

\section{M5：场内期权估值及价差风险模型}
\label{sec:m5-option}

\subsection{模型定位与计量口径}

M5 面向场内期权建立权利金、理论价值、隐含波动率曲面和持有期损失的联合模型。模型首先冻结使用可交易权利金还是无套利理论价值作为主价值口径，并统一合约乘数、看涨/看跌方向、买卖头寸和分红处理。对持有期未到期的期权，损益来自期末权利金或理论价值与期初价值之差；对持有期内到期或行权的合约，则按照最终结算和行权现金流形成 PnL，不能简单删除。

\subsection{经典定价锚、波动率曲面与联合情景}

M5 以 Black--Scholes 或 Black--76 作为产品经典锚\cite{black_scholes_1973,black_1976}。股票期权和 ETF 期权可根据分红与融资口径选用现货框架，期货期权采用与标的期货一致的 Black--76 框架。市场期权链用于反演隐含波动率，并在到期和执行价维度构造满足基本无套利条件的曲面；学习器只拟合曲面动态或经典模型残差，不能用平滑后的事后全样本曲面泄漏未来信息。

模型联合模拟标的价值、利率、分红和隐含波动率曲面，并逐情景重估期权：
\begin{equation}
  V_{t+h}^{(m)}
  =\mathcal{P}\!\left(S_{t+h}^{(m)},K,T-t-h,
  r_{t+h}^{(m)},\sigma_{\mathrm{IV},t+h}^{(m)};\text{terms}\right),
\end{equation}
其中 \(\mathcal{P}\) 为与产品条款相符的定价算子。由价值路径形成权利金或组合 MTM 损失，并输出 VaR、ES、隐含波动率水平与偏斜、Delta、Gamma、Vega、Theta、成交/持仓流动性和对冲误差。临近到期、深度虚实值及低价格样本应采用专门数值规则，避免相对误差失真。

\subsection{基准、验证与证据边界}

M5 的朴素基准为上一期权利金、理论价值或隐含波动率；经典基准为 Black--Scholes/Black--76 与插值隐含波动率曲面；强现代基准为带无套利约束的神经曲面或残差模型。验证需在完全相同的期权链、报价过滤和到期处理规则下，比较价值误差、IV 曲面误差、Greeks 稳定性、对冲误差、分布得分以及 VaR/ES 返回检验。

现有数据包含指数、个股和商品期权的价格、执行价、标的及部分 Greeks，可作为主数据候选；但合约条款统一、无套利过滤、临近到期处理以及 T+10 期间到期合约的可评价性仍需审计。因此 M5 尚不能仅凭部分期权链的拟合结果宣称全品种或持有期风险验证完成。

\section{M6：以场内证券为标的的信用衍生品估值及价差风险模型}
\label{sec:m6-credit-derivative}

\subsection{模型定位与计量口径}

M6 面向以场内证券为标的的信用衍生品，建立合约公允价值、信用利差和违约跳跃损失模型。研究对象必须具有可核验的参考实体或参考证券、保护费支付规则、信用事件定义、结算方式和回收安排，不能把期货或期权数据重新命名为信用衍生品。对保护买方，合约价值原则上表示为保护腿现值减去保费腿现值；保护卖方的符号相反。

持有期总损益由公允价值变化、期间保费、信用事件赔付和结算现金流共同形成：
\begin{equation}
  \operatorname{PnL}_{t,h}
  =V_{t+h}^{\mathrm{fair}}-V_t^{\mathrm{fair}}
  +CF_{t,h}^{\mathrm{premium}}
  +CF_{t,h}^{\mathrm{protection}},
\end{equation}
并据此计算 MTM loss 与 jump-to-default loss。running/par spread、upfront、CS01 和回收率是产品特有指标，但不能替代统一的持有期价值和损失分布。

\subsection{违约强度、生存概率与双腿现金流}

M6 采用 reduced-form 框架描述风险中性违约强度、生存概率和回收率\cite{duffie_singleton_1999}。若 \(\lambda(u)\) 为违约强度，则条件生存概率可写为
\begin{equation}
  Q_t(\tau)=\exp\!\left[-\int_t^{\tau}\lambda(u)\,\mathrm du\right].
\end{equation}
保费腿由未违约状态下的定期费用及应计费用构成，保护腿由信用事件概率、损失率和贴现因子构成。未来情景联合推进基准曲线、信用利差、违约强度、回收率及参考证券状态，在信用事件路径上显式触发赔付和结算，而不是仅对利差做连续扰动。

在真实数据到位后，M6 应输出合约公允价值、未来价值分位数、MTM 损失、jump-to-default、VaR、ES、par/running spread、upfront、CS01、回收率敏感度、流动性和压力损失。基准阶梯依次为上一利差、经典强度/回收率定价以及经验证的面板或时序残差模型；任何强现代方法均须与经典现金流模型共享合约条款和评价样本。

\subsection{数据阻塞与验证边界}

当前项目缺少真实信用衍生品的端到端合约目录、条款、报价、成交、保费腿/保护腿现金流、回收率和信用事件数据，现有“场内证券衍生品”数据实为期权或期货子集，不能作为 M6 的替代输入。因此，M6 当前状态为 \texttt{BLOCKED\_DATA}：系统可以返回带原因码的阻塞证据包，但不得生成看似正常的数值估值或风险结论。

合成数据仅可检验接口、现金流引擎、情景触发和证据归档是否能够运行，结果必须标记为 engineering-only，不得计入真实市场覆盖率、模型准确率或合同通过率。只有取得真实数据，或获得甲方和评审机构对替代范围的书面决定后，才能冻结强现代基准、开展历史返回检验并形成 M6 的正式结论。

